\documentclass[10pt,a4paper]{amsart}
\usepackage[margin=19mm]{geometry}
\usepackage{amsmath,amssymb,mathtools}
\usepackage[scr=rsfs,cal=euler]{mathalfa}
\usepackage{xfrac}
\usepackage[unicode,hidelinks,bookmarks=false]{hyperref}
\newcommand{\ie}{i.e.\ }
\newcommand{\cf}{cf.\ }
\newcommand{\resp}{respectively\ }
\newcommand{\Subsection}{\subsection}
\providecommand{\nobreakdash}{\nobreak\hbox{-}}
\setlength{\emergencystretch}{2em}
\multlinegap0pt
\usepackage{amssymb}
\usepackage{xcolor}
\newtheorem{theorem}{Theorem}
\newtheorem{prop}{Proposition}
\def\Ab{\mathbf{A}}
\def\Bb{\mathbf{B}}
\def\Cb{\mathbf{C}}
\newcommand\RR{\mathbb{R}}
\renewcommand{\ge}{\geqslant}
\renewcommand{\geq}{\geqslant}
\renewcommand{\le}{\leqslant}
\renewcommand{\leq}{\leqslant}
\newcommand\ve{\varepsilon}
\title{Bounds on the exceptional set in the $abc$ conjecture}
\author{Christian Bernert, Tim Browning, Jared Duker Lichtman, Joni Ter\"av\"ainen}
\date{Source: arXiv:2410.12234v2 (CC BY 4.0)}
\begin{document}
\maketitle

\noindent\emph{Source and licence.} Bounds on the exceptional set in the $abc$ conjecture. Version arXiv:2410.12234v2 (CC BY 4.0), distributed by arXiv under the Creative Commons Attribution 4.0 International licence, http://creativecommons.org/licenses/by/4.0/, as recorded in the arXiv licence metadata for this exact version and shown on the abstract page https://arxiv.org/abs/2410.12234v2. Full licence notice: \texttt{licenses/arxiv-2410.12234-CC-BY-4.0.txt}.\par\smallskip

\noindent\textit{Editorial scope.} Counted unit: the article's theorem thm:main\_optimal (source0.tex lines 197--199). The article's complete original proof of it is delivered with it (source0.tex lines 597--650, one \texttt{proof} environment of 54 lines, which the article places away from the statement). No mathematics is written, repaired or re-derived: the statement and the proof are the article's own lines, in the article's own notation, and no retained line is substituted or re-flowed.  Retained uncounted as the source-local context the proof needs: lines 268--279; lines 366--369; lines 378--381; lines 426--431; lines 490--593; lines 496--498; lines 502--510; lines 514--522; lines 526--528; lines 532--534; lines 538--546; lines 550--558; lines 562--570; lines 574--582; lines 586--592.  Nothing was omitted from any retained span, so the package carries no omission record.  No retained line contains a citation, so the wrapper carries no bibliography and no bibliography entry.  No retained line contains a figure environment, an image-inclusion command or drawing code, so no image is delivered and the package declares no asset dependency.  Editorial formatting only, in addition: an AMS \texttt{amsart} preamble that restates the article's own declarations and macros used by the retained lines, namely the environments \texttt{theorem}, \texttt{prop} on separate counters, together with the standard packages those lines need.  The retained environments print in this document as Theorem 1, Proposition 1 in order of appearance, whereas the article prints the same items with the numbering of its own full document.  The complete native source of the article as distributed in the arXiv e-print for this exact version is bundled unchanged as \texttt{sources/166912/source0.tex}.
\begin{prop}\label{prop:reduction}
    For any $\varepsilon_0>0$, there exists a positive integer $d=d(\varepsilon_0)$ such that
    \[
    N_{\lambda}^*(X) \ll_{\varepsilon_0} X^{\varepsilon_0} \cdot \max B_d(a_0,b_0,c_0,\Ab,\Bb,\Cb),
    \]
    where the maximum runs over pairwise coprime $1\le a_0,b_0,c_0 \le X^{\varepsilon_0}$ and dyadic parameters $A_i,B_i,C_i \in [1,X]$ with
    \[
    \frac{1}{2^{d+1}}X^{\lambda}\leq 
    \prod_{j=1}^d A_jB_jC_j \le X^{\lambda+3\varepsilon_0/4}\]
    and
    \[\prod_{j=1}^d A_j^j \le X, \quad \prod_{j=1}^d B_j^j \le X, \quad X^{1-\varepsilon_0/2} \le \prod_{j=1}^d C_j^j \le X.\]
\end{prop}

\begin{prop}\label{prop:triv2}
    We have
    \[B_d(a_0,b_0,c_0,\Ab,\Bb,\Cb) \ll_{\varepsilon} X^{\varepsilon}\min(AB,AC,BC).\]
\end{prop}

\begin{prop}\label{prop:fourierbound}
    For $2 \le j,k,\ell \le d$, we have
    \[B_d(a_0,b_0,c_0,\Ab,\Bb,\Cb) \ll_{\varepsilon_0} \frac{X^{2\lambda/3+3\varepsilon_0}}{\left(\prod_{j \mid i} A_i \prod_{k \mid i} B_i \prod_{\ell \mid i} C_i\right)^{1/6}}.\]
\end{prop}

\begin{prop}\label{prop:geometrybound}
    For any sets $I,I',I'' \subset \{1,2,\dots,d\}$ we have
    \begin{align*}
        B_d(a_0,b_0,c_0,\Ab,\Bb,\Cb) \ll_{\varepsilon_0}~& \frac{X^{\lambda+4\varepsilon_0}}{\prod_{i \in I} A_i \prod_{i \in I'} B_i \prod_{i \in I''} C_i} \\
        &+X^{\lambda-1+6\varepsilon_0} \prod_{i \in I} A_i^{i-1} \prod_{i \in I'} B_i^{i-1} \prod_{i \in I''} C_i^{i-1}.\end{align*}
\end{prop}

\begin{prop}\label{prop:linear_program}
    Let $a,a_1,\dots,a_6,b,b_1,\dots,b_6,c,c_1,\dots,c_6,D \in \RR_{\geq 0}$. Let  
    $$
    p_i=a_i+b_i+c_i, \quad L=a+b+c. 
     $$
 Then $D \le 0.6$ if the  following list of $22$ inequalities hold:
    \begin{equation}\label{eq:1}
            L \le 1,
    \end{equation}

\medskip

    \begin{equation}\label{eq:2}
\left\{
\begin{aligned}
a_1+a_2+a_3+a_4+a_5+a_6 &\le a,\\
b_1+b_2+b_3+b_4+b_5+b_6 &\le b,\\
c_1+c_2+c_3+c_4+c_5+c_6 &\le c,
\end{aligned}
\right.
\end{equation}

\medskip

\begin{equation}\label{eq:3}
\left\{
\begin{aligned}
        7a-1 &\le 6a_1+5a_2+4a_3+3a_4+2a_5+a_6,\\
        7b-1 &\le 6b_1+5b_2+4b_3+3b_4+2b_5+b_6,\\
        7c-1 &\le 6c_1+5c_2+4c_3+3c_4+2c_5+c_6,
\end{aligned}
\right.
\end{equation}

\medskip

        \begin{equation}\label{eq:4}
        D \le \min(a+b,a+c,b+c),
    \end{equation}

\medskip

\begin{equation}\label{eq:5}
       \min(p_1+p_2+p_4,1-p_2-3p_4) \le L-D,
    \end{equation}

\medskip

\begin{equation}\label{eq:6}
\left\{
\begin{aligned}
        \min(p_1+p_2+a_3,1-p_2-2a_3) &\le L-D,\\
        \min(p_1+p_2+b_3,1-p_2-2b_3) &\le L-D,\\
        \min(p_1+p_2+c_3,1-p_2-2c_3) &\le L-D,
\end{aligned}
\right.
\end{equation}

\medskip

\begin{equation}\label{eq:7}
\left\{
\begin{aligned}
        \min(p_1+p_2+a_3+b_3,1-p_2-2a_3-2b_3) &\le L-D,\\
        \min(p_1+p_2+b_3+c_3,1-p_2-2b_3-2c_3) &\le L-D,\\
        \min(p_1+p_2+c_3+a_3,1-p_2-2c_3-2a_3) &\le L-D,
\end{aligned}
\right.
\end{equation}

\medskip

\begin{equation}\label{eq:8}
\left\{
\begin{aligned}
        \min(p_1+a_3+b_2+c_2,1-2a_3-b_2-c_2) &\le L-D,\\
        \min(p_1+a_2+b_3+c_2,1-a_2-2b_3-c_2) &\le L-D,\\
        \min(p_1+a_2+b_2+c_3,1-a_2-b_2-2c_3) &\le L-D,\\
\end{aligned}
\right.
\end{equation}

\medskip

\begin{equation}\label{eq:9}
\left\{
\begin{aligned}
        \min(p_1+a_2+b_3+c_3,1-a_2-2b_3-2c_3) &\le L-D,\\
        \min(p_1+a_3+b_2+c_3,1-2a_3-b_2-2c_3) &\le L-D,\\
        \min(p_1+a_3+b_3+c_2,1-2a_3-2b_3-c_2) &\le L-D,\\
\end{aligned}
\right.
\end{equation}

\medskip

\begin{equation}\label{eq:10}
\begin{split}
        \max(a_2+a_4+a_6,a_3+a_6,a_5)&+\max(b_2+b_4+b_6,b_3+b_6,b_5)\\
        &+\max(c_2+c_4+c_6,c_3+c_6,c_5) 
        \le 4L-6D.
\end{split}
\end{equation}
\end{prop}

    \begin{equation}
            L \le 1,
    \end{equation}

    \begin{equation}
\left\{
\begin{aligned}
a_1+a_2+a_3+a_4+a_5+a_6 &\le a,\\
b_1+b_2+b_3+b_4+b_5+b_6 &\le b,\\
c_1+c_2+c_3+c_4+c_5+c_6 &\le c,
\end{aligned}
\right.
\end{equation}

\begin{equation}
\left\{
\begin{aligned}
        7a-1 &\le 6a_1+5a_2+4a_3+3a_4+2a_5+a_6,\\
        7b-1 &\le 6b_1+5b_2+4b_3+3b_4+2b_5+b_6,\\
        7c-1 &\le 6c_1+5c_2+4c_3+3c_4+2c_5+c_6,
\end{aligned}
\right.
\end{equation}

        \begin{equation}
        D \le \min(a+b,a+c,b+c),
    \end{equation}

\begin{equation}
       \min(p_1+p_2+p_4,1-p_2-3p_4) \le L-D,
    \end{equation}

\begin{equation}
\left\{
\begin{aligned}
        \min(p_1+p_2+a_3,1-p_2-2a_3) &\le L-D,\\
        \min(p_1+p_2+b_3,1-p_2-2b_3) &\le L-D,\\
        \min(p_1+p_2+c_3,1-p_2-2c_3) &\le L-D,
\end{aligned}
\right.
\end{equation}

\begin{equation}
\left\{
\begin{aligned}
        \min(p_1+p_2+a_3+b_3,1-p_2-2a_3-2b_3) &\le L-D,\\
        \min(p_1+p_2+b_3+c_3,1-p_2-2b_3-2c_3) &\le L-D,\\
        \min(p_1+p_2+c_3+a_3,1-p_2-2c_3-2a_3) &\le L-D,
\end{aligned}
\right.
\end{equation}

\begin{equation}
\left\{
\begin{aligned}
        \min(p_1+a_3+b_2+c_2,1-2a_3-b_2-c_2) &\le L-D,\\
        \min(p_1+a_2+b_3+c_2,1-a_2-2b_3-c_2) &\le L-D,\\
        \min(p_1+a_2+b_2+c_3,1-a_2-b_2-2c_3) &\le L-D,\\
\end{aligned}
\right.
\end{equation}

\begin{equation}
\left\{
\begin{aligned}
        \min(p_1+a_2+b_3+c_3,1-a_2-2b_3-2c_3) &\le L-D,\\
        \min(p_1+a_3+b_2+c_3,1-2a_3-b_2-2c_3) &\le L-D,\\
        \min(p_1+a_3+b_3+c_2,1-2a_3-2b_3-c_2) &\le L-D,\\
\end{aligned}
\right.
\end{equation}

\begin{equation}
\begin{split}
        \max(a_2+a_4+a_6,a_3+a_6,a_5)&+\max(b_2+b_4+b_6,b_3+b_6,b_5)\\
        &+\max(c_2+c_4+c_6,c_3+c_6,c_5) 
        \le 4L-6D.
\end{split}
\end{equation}

\begin{theorem} \label{thm:main_optimal}
Let $\lambda \in (0,1)$. Then $N_\lambda(X) \ll_{\ve,\lambda} X^{0.6+\ve}$, for any $\ve>0$.
\end{theorem}

\begin{proof}[Proof of Theorem~\ref{thm:main_optimal}]
We fix a choice of  $\lambda\in (0,1)$. 
Suppose that 
Theorem~\ref{thm:main_optimal} is false. Then, by Proposition~\ref{prop:reduction}, we can find $\delta_0>0.6$ such that for any $\varepsilon_0$ and the corresponding $d$, there are $a_0,b_0,c_0, \Ab,\Bb,\Cb$ with
    \[B_d(a_0,b_0,c_0,\Ab,\Bb,\Cb) \ge X^{\delta_0}\]
    for arbitrarily large $X$. Fix $D \in (0.6,\delta_0)$.

    Write $A=X^a, B=X^b, C=X^c$ as well as $A_i=X^{a_i}, B_i=X^{b_i}, C_i=X^{c_i}$. Then, as we will see below, for sufficiently small $\varepsilon_0>0$, the variables $a_i,b_i,c_i,a,b,c,D$ satisfy the hypotheses of Proposition~\ref{prop:linear_program} contradicting the fact that $D>0.6$.

Assuming that $\ve_0$ is sufficiently small, it follows that $ABC\leq X^{\lambda+3\ve_0}\leq X$. This establishes  \eqref{eq:1}. The 
inequalities in
  \eqref{eq:2} 
follow immediately from the definitions. 
The inequalities in \eqref{eq:3} also follow from the definitions and the observation that 
$$
\sum_{1\leq i\leq 6} ia_i \leq 1, \quad 
\sum_{1\leq i\leq 6} ib_i \leq 1, \quad
\sum_{1\leq i\leq 6} ic_i \leq 1,
$$
so that 
$$
7a-1\le 7a-\sum_{1\leq i\leq 6} ia_i\leq 6a_1+5a_2+4a_3+3a_4+2a_5+a_6.
$$
Similarly for the remaining two inequalities in \eqref{eq:3}.
The  inequality \eqref{eq:4} is just the trivial bound from Proposition~\ref{prop:triv2}.


    The next set of  inequalities originate from the geometry bound in Proposition~\ref{prop:geometrybound} with various choices of $I$, $I'$, $I''$. 
       More precisely, \eqref{eq:5} follows on taking
    \[
    I=I'=I''=\{1,2,4\};
    \]
the inequalities \eqref{eq:6} correspond to taking 
    \[
    I=I'=\{1,2\}, I''=\{1,2,3\}
    \]
    and its permutations;
the inequalities \eqref{eq:7} correspond to taking 
    \[
    I=\{1,2\}, I'=I''=\{1,2,3\}
    \]
    and its permutations; 
    the inequalities \eqref{eq:8} correspond to taking 
    \[
    I=I'=\{1,2\}, I''=\{1,3\}
    \]
    and its permutations; and finally
the inequalities \eqref{eq:9} correspond to taking 
    \[
    I=\{1,2\}, I'=I''=\{1,3\}
    \]
    and its permutations.
The final  inequality \eqref{eq:10} is the Fourier bound from Proposition~\ref{prop:fourierbound}.
\end{proof}

\end{document}
